\section{A finite geometric test for disc embedding}\label{sec:geometry}\label{geo:section}

We construct the input discs and their algebraic duals in an actual smooth
four-manifold.  The geometric conclusion needed later is conditional: any
system of disjoint proper topological discs with the prescribed boundaries
produces the smooth cut data specified in Proposition~\ref{geo:construction}.
In particular, the reduction places no relative homotopy requirement on
those hypothetical discs.  Intersection and reduced self-intersection
invariants have the conventions of \cite[Section~1.3]{PRT2025}; no disc-embedding
theorem is used in this section.

\subsection{The ambient manifold and its boundary}

Fix an integer $d\geq5$.  For each $e\in\{1,\ldots,d\}$ take two
end labels $(e,0),(e,1)$, and put
\[
 \mathcal E=\{(e,\varepsilon):1\leq e\leq d,
                          \ \varepsilon\in\{0,1\}\}.
\]
The two ends have no preferred positive or negative designation yet.

Let $V$ be the oriented four-dimensional $1$-handlebody with $3+2d$
handles.  Label $2d$ of its handle generators by the ends, and name the
remaining three $h_0,h_+,h_-$.  Its boundary is a connected sum of
$3+2d$ copies of $S^1\times S^2$.  In a separate punctured summand for each end $a$,
choose a standard generator knot $K_a=S^1\times\{\mathrm{point}\}$
and a closed tubular solid torus $N_a$.  Choose all connected-sum balls
away from these tori.  On $\partial N_a$ write $\ell_a$ for the
product longitude and $m_a$ for the meridian.  Define
\[
 Q=\overline{\partial V\setminus\bigcup_{a\in\mathcal E}
                                      \operatorname{int}N_a}.
\]
The choices of basepoint and paths in $Q$ will be fixed throughout.

For each $e$, attach an oriented $4$-ball $P_e$ to $V$ along $N_{e,0}$
and $N_{e,1}$.  In $\partial P_e=S^3$ these are tubular neighborhoods
of the two components of the standard Hopf link.  Choose the gluing
maps to identify the generator and product framing with those on
$N_{e,0}$ and $N_{e,1}$, using the signs that give an oriented gluing.
Equivalently, take the rounded product ball $D^2\times D^2$ and its two
coordinate core discs.  Denote the resulting compact smooth manifold,
after rounding corners, by $G$.  The remaining boundary portion of
$P_e$ is a Hopf-link exterior
\begin{equation}\label{geo:hopf-collar}
 C_e\cong T^2\times[0,1].
\end{equation}
Under its two end identifications, $\ell_{e,0}$ is $m_{e,1}$ and
$m_{e,0}$ is $\ell_{e,1}$, up to orientation signs.  Thus the two
longitude slopes become distinct coordinate circles on one common
torus.  The two core discs are framed and cross once; their small
parallels have the explicit form
\begin{equation}\label{geo:product-sheets}
 D^2\times\{a_\nu\},\qquad \{b_\mu\}\times D^2,
\end{equation}
with distinct parameters in each family.  Each opposite-family pair
crosses once and each same-family pair is disjoint.

\begin{lemma}[Graph model and peripheral group]\label{geo:graph}
The manifold $Q$ is connected and
\[
 \pi_1(Q)=F(h_0,h_+,h_-,z_a\ (a\in\mathcal E)),
 \qquad [\ell_a]=z_a,\quad [m_a]=1.
\]
The manifold $G$ has the homotopy type of the rose
$\Gamma$ on $h_0,h_+,h_-,q_1,\ldots,q_d$.  There is a homotopy
equivalence $p_0:G\longrightarrow\Gamma$ that sends $V$ into the
$h$-subrose $\Gamma_h$, induces the named generators on $h$, kills
the end generators, and on each $C_e$ is its interval coordinate
traversing $q_e$ once.
\end{lemma}

\begin{proof}
In an end summand, deleting $S^1\times\operatorname{int}D^2$
leaves $S^1\times D^2$ with a ball removed.  The old meridian bounds
the complementary $D^2$, whereas the old longitude generates its
fundamental group.  Van Kampen across the connected-sum spheres gives
the displayed free product and proves connectedness.

All attachments are collared cofibrations, so the ordinary gluing
computes the homotopy pushout.  Replace $V$ by its bouquet of handle
circles, each $N_a$ by its generator circle, and each $P_e$ by a point.
For an edge $e$ the resulting pushout cones off the two independent
circles $z_{e,0},z_{e,1}$ to the same new point.  Each of the two cone
discs retracts onto an arc from the old bouquet point to this cone
point, keeping both endpoints fixed.  The two arcs together give one
circle $q_e$.  The other circles are the three $h$ circles.  The only point of each
boundary circle shared with the other summands is the old bouquet point, so these retractions are compatible.
This proves the asserted graph homotopy type.

On $V$ take the handle map retaining the $h$ circles and constant on
each $N_a$.  On $P_e$ prescribe a function to $[0,1]$ equal to $0$ and
$1$ on its two attaching tori and to the product coordinate on $C_e$.
The compatible boundary function extends over the ball: extend it as
a real function and compose with the retraction of $\mathbb R$ onto
$[0,1]$.  Compose with the edge $q_e$, whose endpoints are the common
vertex.  The resulting map induces the graph identification just
described, and is therefore a homotopy equivalence.  All boundary
values may be taken productwise in collars.
\end{proof}

Figure~\ref{geo:fig-graph} records both parts of this construction: the
new graph loop and the interchange of the two peripheral slopes.

\begin{figure}[!hb]
\centering
\begin{tikzpicture}[font=\small,>=Stealth,
  block/.style={draw,rounded corners,align=center,inner sep=6pt},
  every path/.style={line width=.5pt}]
  \node[block,minimum width=2.7cm,minimum height=1.7cm] (V) at (0,0)
    {$V$\\$h,z_{e,0},z_{e,1},\ldots$};
  \node[block,minimum width=2.4cm,minimum height=1.7cm] (P) at (4.3,0)
    {$P_e=D^2\times D^2$\\two core cap types};
  \draw (V.north east) -- node[above,sloped] {$N_{e,0}$} (P.north west);
  \draw (V.south east) -- node[below,sloped] {$N_{e,1}$} (P.south west);
  \node[align=center] at (2.1,-1.7)
    {Remaining boundary: $C_e=T^2\times[0,1]$\\
     $\ell_{e,0}\leftrightarrow m_{e,1}$,\quad
     $m_{e,0}\leftrightarrow\ell_{e,1}$};
  \draw[->] (6.05,0) -- (6.8,0);
  \node[above,align=center] at (6.4,.65) {homotopy\\model};
  \fill (7.8,0) circle (2pt);
  \fill (9.8,0) circle (2pt);
  \draw (7.8,0) .. controls (8.4,1) and (9.2,1) .. (9.8,0);
  \draw (7.8,0) .. controls (8.4,-1) and (9.2,-1) .. (9.8,0);
  \draw (7.8,0) .. controls (6.9,1.3) and (6.9,-1.3) .. (7.8,0);
  \node at (7.15,-.75) {$h$};
  \node at (8.8,0) {$q_e$};
  \node[align=center] at (8.7,-1.7)
    {Two end circles contract;\\two remaining arcs give one loop.};
\end{tikzpicture}

\caption{Incidence and peripheral data for one special edge.  This is
a gluing schematic, not a drawing of an embedding of the four-manifold
in three-dimensional space.  The Hopf exterior interchanges longitude
and meridian; the ball kills two end generators and contributes the
single graph loop $q_e$.}\label{geo:fig-graph}
\end{figure}

\subsection{The word list and capped stages}

Write $[u,v]=uvu^{-1}v^{-1}$ and $z^c=czc^{-1}$.  Let $F_h$ be
free on the three named generators $h_0,h_+,h_-$.  Take $\mathcal H$
to be all reduced words of length at most five in
$h_0^{\pm1},h_+^{\pm1},h_-^{\pm1}$, including the empty word.
This is a finite set closed under inversion and reversal.  Put
\[
                  \mathcal R=(\mathcal E\times\mathcal H)^3.
\]
Thus $k=|\mathcal R|=(2d\,|\mathcal H|)^3$ is finite.  The entry
$r=((a,c_1),(b,c_2),(c,c_3))$ specifies the word
\begin{equation}\label{geo:relations}
 R_r=[z_a^{c_1},[z_b^{c_2},z_c^{c_3}]].
\end{equation}
All triples occur, including repetitions, and the three conjugators vary
independently.

We realize each $R_r$ by two punctured tori.  The base $B_r$ has
basis curves $\alpha_r,\beta_r$: the first receives a cap of end type
$a$, while the second receives a punctured torus $S_r$ for the inner
commutator.  Its basis curves $x_r,y_r$ receive caps of end types $b,c$.
Compressing $S_r$ with two copies of its $x_r$ cap will provide the
contracting disc used to construct a dual sphere.  We retain the unused
$y_r$ cap as well; it will identify the group-ring labels of paired
intersections in the exterior of the bases.

A \emph{body} means one of these embedded punctured tori, excluding its
caps.  Bodies and their attaching collars will be disjoint except at
the prescribed attachments.  Opposite end cap types are allowed to
cross in $P_e$.  We first explain how the ribbon construction controls
the attachment framings, then assemble the finite word list.

\begin{lemma}[The framed ribbon operation]\label{geo:framing}
Given two disjoint oriented framed knots in the interior of $Q$ and a
relative path class joining them, one can construct a punctured torus
in a level of a collar $Q\times[0,\epsilon)$ whose two basis curves
resolve normally to these knots.  Its oriented boundary represents
their commutator, with the change of basing given by the chosen path.
The two band framings can be varied by arbitrary integers independently.
The resolved knots admit disjoint rising annuli to the given framed
knots, and a unit band twist changes the relative oriented normal
framing of its annulus by a unit integer.
\end{lemma}

\begin{proof}
Represent the path by an embedded arc disjoint from the knots except
at its endpoints.  Along a narrow finger following this arc move a
short strand of the second knot to the first, stopping when the two
strands meet transversely in a small two-dimensional patch.  Before
the final contact the move is an isotopy.  Thicken the resulting
crossed graph to an oriented ribbon with alternating cyclic order
of the four half-edges at the crossing.  Its regular neighborhood is
a punctured torus.  Choose the normal resolution determined by the
pre-contact strands.  Reversing that specific finger isotopy recovers
the original two-component link.  The boundary traverses the circuits
in the order
$uvu^{-1}v^{-1}$, after a choice of orientation and basing.  The two
sides of the finger insert precisely the chosen change-of-basing
path and its inverse.

The ribbon away from the crossing consists of separate bands.
Insert full twists in either band, supported in a short segment away
from the crossing and the other band.  Take its supporting tube disjoint
from the other resolved knot as well.  Shrinking the displacement in this tube gives
an ambient isotopy of the resolved link that fixes its other component.
Thus a twist preserves the chosen two-component resolution, including
its linking.  The lateral ribbon vector rotates once per full twist,
giving independent integer changes to the two knot framings.

Here is the normal-bundle calculation needed to transfer that integer
to a cap attachment.  Let $t$ be tangent to a basis curve, $b$ its
lateral ribbon vector, $n$ the normal to the ribbon in the level
three-manifold, and $v$ the positive collar direction.  A rising
annulus initially has tangent plane spanned by $t$ and
$r=a n+c v$, with $a^2+c^2=1$ and $c>0$.  Set $w=c n-a v$.
Its oriented normal quotient has frame $([b],[w])$.  Subtracting
the $v$ component divided by $c$ times $r$ identifies this quotient
with the normal plane of the movie knot in the level slice, and gives
\begin{equation}\label{geo:normal-quotient}
 [b]\longmapsto[b],\qquad [w]\longmapsto[n/c].
\end{equation}
Thus it carries exactly the ribbon framing, with a positive rescaling
of the second vector.  During one band twist the pair changes to
\[
 b_\theta=\cos\theta\,b+\sin\theta\,n,
 \qquad n_\theta=-\sin\theta\,b+\cos\theta\,n,
\]
where $\theta$ makes a full turn.  Equation~\eqref{geo:normal-quotient}
has degree one in the oriented rank-two frame group, up to the chosen
sign convention.  This is an integer calculation in
$\pi_1(\operatorname{GL}^+(2,\mathbb R))$, rather than a calculation
of the parity of an ambient three-frame.

The separation of the attaching annuli can be checked in local collar
coordinates $(x,y,n,v)$, with the body in $n=v=0$ and its basis curves
the $x$- and $y$-axes.  Choose the signs of these coordinates to match
the selected normal resolution, and take the initial annuli to be
\[
 (x,0,s,s),\qquad (0,y,-s,s),\qquad s\geq0.
\]
They meet on the body at $s=0$ and are disjoint for $s>0$.  At a small
positive level they give the chosen normal resolution; continue them
by the inverse finger isotopy of the whole link.  The band-twist
isotopies just described are supported away from the other component,
so their traces also preserve this disjointness.
Figure~\ref{geo:fig-ribbon} separates the body coordinates from the
normal coordinates in this local model.  Along a rising movie the same
quotient identification transports the framing to the given knot.
Consequently a full band twist changes by one the obstruction to
extending a prescribed ordered normal frame over an upper cap disc.
Killing that integer extends both vectors: an extended nonzero first
vector and its oriented quarter-turn give a frame, whose second
vector can be adjusted at the boundary through positive independent
vectors.  This also explains why the statement controls a prescribed
boundary framing, not just abstract triviality of a disc bundle.
\end{proof}

\begin{figure}[tb]
\centering
\begin{tikzpicture}[font=\small,>=Stealth,
  every path/.style={line width=.6pt}]
  \begin{scope}[shift={(0,0)}]
    \fill[gray!10] (-1.9,-1.25) rectangle (1.9,1.25);
    \draw[->] (-1.7,0) -- (1.7,0) node[right] {$x$};
    \draw[->] (0,-1.05) -- (0,1.05) node[above] {$y$};
    \fill (0,0) circle (2pt);
    \node[align=center] at (0,-1.8)
      {Body plane $n=v=0$\\basis curves meet once};
  \end{scope}
  \begin{scope}[shift={(6,0)}]
    \draw[->,gray] (-2,0) -- (2,0) node[right] {$n$};
    \draw[->,gray] (0,-.2) -- (0,1.65) node[above] {$v$};
    \draw[->,thick] (0,0) -- (1.3,1.3)
      node[right,align=left] {$x$-annulus\\$n=v=s$};
    \draw[->,thick] (0,0) -- (-1.3,1.3)
      node[left,align=right] {$y$-annulus\\$n=-s,\ v=s$};
    \draw[dashed] (-1.6,.75) -- (1.6,.75);
    \fill (-.75,.75) circle (2pt);
    \fill (.75,.75) circle (2pt);
    \node[anchor=west,inner sep=2pt] at (1.65,.75) {$v=s>0$};
    \node[align=center] at (0,-1.8)
      {Normal plane $(n,v)$\\separate points at each positive level};
  \end{scope}
\end{tikzpicture}

\caption{Local coordinates for the ribbon attachment.  The left panel
shows the intersecting basis curves in the body plane; the right panel
shows their distinct rising rays in the normal $(n,v)$-plane.  At each
positive collar level the two annuli are separated in $n$.  The chosen
resolution is the one recovered by reversing the pre-contact finger
isotopy.  This local picture makes no assertion that cap images are
mutually disjoint.}\label{geo:fig-ribbon}
\end{figure}

\begin{lemma}[Simultaneous capped stages]\label{geo:stages}
For every $r\in\mathcal R$ there is a punctured torus $B_r$ properly
embedded in $V$, with boundary a knot in $\operatorname{int}Q$.
The $B_r$ are mutually disjoint.  They have intersecting basis curves
$\alpha_r,\beta_r$ and the following attachments in $G$:
\begin{enumerate}
\item A shallow cap at $\alpha_r$ ends in a separate parallel copy of
the core disc of end $a$.
\item A further punctured torus $S_r$ is attached at $\beta_r$.
Its basis curves $x_r,y_r$ have caps of end types $b,c$, respectively.
Its boundary, with the prescribed basing, represents
$[z_b^{c_2},z_c^{c_3}]$, and $\partial B_r$ represents $R_r$.
\item Compressing $S_r$ along $x_r$ with two framed parallels of its
cap produces an embedded disc $D_r$.  Intersections of $D_r$ with
other objects are allowed.  The extension of $D_r$ towards $\beta_r$
has a framing that extends the angular normal-circle vector required
to compress $\beta_r\times S^1$.
\end{enumerate}
Except at the intended attachments, every body and attaching collar
is disjoint from every other such piece.  All other intersections
are transverse crossings of opposite end cap sheets in the product
charts \eqref{geo:product-sheets}.
\end{lemma}

\begin{proof}
Give each leaf occurrence its own parallel longitude knot on the
$Q$ side of its end torus, even when several occurrences have the
same label.  Reserve a separate return lane inside the corresponding
solid torus.  Each lane leads to its own framed core-disc parallel
in $P_e$.  There are finitely many occurrences, so all lanes may be
chosen disjoint within each end family.

Choose narrow, separated ranges of collar levels centered at
\[
                    0<t_B<t_S<t_T<\epsilon.
\]
The base-body operations occupy the range about $t_B$, the inner-body
operations the range about $t_S$, and the transfers into the end solid
tori the range about $t_T$.  Each body-operation range includes room
above its body level for the rising resolution and inverse finger isotopy of
Lemma~\ref{geo:framing}.  Between ranges, continuing knots run in narrow
disjoint pipes; at an active range all inactive continuations remain in
their pipes.

To trace one relation through these heights, its base boundary rises
from level zero to $B_r$ near $t_B$.  The shallow attachment at
$\alpha_r$ continues to the transfer range.  The other attachment rises
from $\beta_r$ to the boundary of $S_r$ near $t_S$; the two attachments
at $x_r,y_r$ then rise through their own pipes to the transfer range.
There every leaf moves from the $Q$ side into its end solid torus and
returns down its reserved lane in the torus interior to its cap in $P_e$.
These downward lanes lie outside the $Q$-side body neighborhoods.
The construction order is upper stage first: after making $S_r$, we
carry its boundary down to a handoff knot at the upper side of the
base-operation range.  With this knot and the slot-$1$ leaf fixed above
the lower body level, we choose their connecting arc.  The lower
ribbon's rising movie then joins the body to those two inputs.  Thus the
later base operation is below the earlier upper body and avoids the
inactive leaf pipes.

For either body-operation range, once its input knots have been fixed,
choose the connecting arcs with separate endpoints and the prescribed relative
path classes in $Q$.  General position in this three-manifold makes
the finitely many arcs embedded, mutually disjoint and disjoint from
the other knots.  No prescribed path class in a fixed link complement
is needed.  Small neighborhoods of these graphs accommodate the
ribbons, isotopies and local twists.

Choose the connecting arcs for the slot-$2,3$ pairs as disjoint arcs
in $Q$, and first apply Lemma~\ref{geo:framing} to every pair of leaf
occurrences in slots $2,3$, to make $S_r$, avoiding all slot-$1$ leaf knots.
Attach the rising annuli to the leaf knots, and match the $x_r$ band
framing to the framed $x_r$ cap.  Cutting out an annulus along $x_r$
leaves a pair of pants.  Capping its two new boundary circles by
nearby copies of the $x_r$ cap gives a disc $D_r$.  The extended
lateral normal field supplied by Lemma~\ref{geo:framing} places
the two copies at the cut band sides.  This disc is embedded:
the retained body and its collars are disjoint, and the only cap
pieces retained have the same end type.  In particular, crossings
with the unused $y_r$ cap or with another leaf are not self-crossings
of $D_r$.

Return the boundary knot of $S_r$ along an attaching collar to a
lower level.  Its contracting disc supplies a normal-framing class
at this handoff.  With the returned boundary knots now fixed, choose
the lower connecting arcs and pair each knot with its slot-$1$ leaf, again by
Lemma~\ref{geo:framing}, to make $B_r$.  Choose the
$\beta_r$ band twist so that the transported ordered frame extends
over this already constructed contracting disc.  At the chosen
normal phase of $\beta_r\times S^1$, the vector $w$ in
\eqref{geo:normal-quotient} is its angular tangent.  Extension of
the ordered frame therefore extends this particular angular
vector, as required for compression of the product torus.

These choices are sequential.  If $n_x,n_\beta$ are the two twists,
the two relative framing obstructions have the form
\[
 o_x(n_x,n_\beta)=o_x(0,0)+\varepsilon_x n_x,
 \qquad
 o_\beta(n_x,n_\beta)=o_\beta(n_x,0)
                           +\varepsilon_\beta n_\beta,
 \qquad \varepsilon_x,\varepsilon_\beta\in\{1,-1\}.
\]
The lower $\beta_r$ twist does not change the fixed upper disc or
its $x_r$ framing.  One first sets $o_x=0$, records the resulting
handoff framing, and then sets $o_\beta=0$.  There is no simultaneous
parity assumption hidden in these choices.

Chosen changes of basing make the leaf circuits the conjugates in
\eqref{geo:relations}.  Applying the commutator rule first to $S_r$
and then to $B_r$ gives the asserted words.  Paths recording basings
need not themselves be physical attachments from a single basepoint.

Finally extend the base boundaries down to level zero in $Q$.
This makes the $B_r$ proper in $V$.  Choose the remaining small
parallel displacements and radii after all finitely many bodies
and collars have been fixed.  Outside the attachment charts and
the cap-crossing charts their compact disjoint pieces have positive
separation; inside the charts use the product choices
\eqref{geo:product-sheets}.  This gives simultaneous smooth choices
with exactly the asserted intersections.
\end{proof}

\subsection{The exterior, the dual spheres, and their labels}

Delete small open relative tubular neighborhoods of the bases only:
\begin{equation}\label{geo:exterior}
 M=\overline{G\setminus\bigcup_{r\in\mathcal R}
                                      \nu(B_r)}.
\end{equation}
The relative neighborhoods meet $\partial G$ in the corresponding
boundary-knot neighborhoods.  Choose their radii small enough to
miss every upper-stage piece except its prescribed initial
attachment.  Round corners with compatible collars.  This is a
compact oriented smooth four-manifold with nonempty boundary.
It is connected: paths may be moved off finitely many embedded
surfaces in dimension four, and the complement of the surfaces
retracts radially onto the exterior of sufficiently small tubes.

Truncate each shallow cap at the boundary of the corresponding
base tube and call the resulting proper disc $f_r$.  Its boundary
circles are mutually disjoint smooth embeddings in $\partial M$.
The cap construction gives smooth disc maps, made generic while
keeping the product collars fixed.  In fact each individual disc
is embedded; discs of opposite end types may meet one another.
Their normal bundles supply their induced boundary framings.

At a slightly larger radius than the deleted tube, put
\[
 T_r=\beta_r\times S^1\ \subset M.
\]
This product torus meets the shallow attaching annulus at exactly
one transverse point, since $\alpha_r$ and $\beta_r$ cross once.
At the distinct phase used by the further stage, compress $T_r$
along its $\beta_r$ direction using two parallels of the extended
contracting disc $D_r$.  Denote the resulting sphere by $g_r$.
Figure~\ref{geo:fig-stages} distinguishes the two successive pairs of
cap copies in this sphere.

\suppressfloats[t]
\begin{figure}[!ht]
\centering
\begin{tikzpicture}[font=\small,>=Stealth,
  body/.style={draw,rounded corners,align=center,inner sep=5pt},
  capnode/.style={draw,rounded corners=8pt,align=center,inner sep=4pt},
  every path/.style={line width=.5pt}]
  \node[body] (base) at (0,0) {$B_r$\\$\partial B_r=R_r$};
  \node[capnode] (shallow) at (-1.65,1.8) {shallow cap\\$z_a^{c_1}$};
  \node[body] (upper) at (1.55,1.8) {$S_r$\\$[z_b^{c_2},z_c^{c_3}]$};
  \node[capnode] (x) at (.3,3.6) {$x_r$ cap\\$z_b^{c_2}$};
  \node[capnode] (y) at (2.9,3.6) {$y_r$ cap\\$z_c^{c_3}$};
  \draw (base) -- node[pos=.33,left] {$\alpha_r$} (shallow);
  \draw (base) -- node[pos=.33,right] {$\beta_r$} (upper);
  \draw (upper) -- node[left] {$x_r$} (x);
  \draw (upper) -- node[right] {$y_r$} (y);
  \draw[dashed] (-2.65,.8) -- (3.6,.8);
  \node[align=center] at (.4,-.8) {Only the base tube is deleted.};

  \node[align=center] at (7,4.15)
    {$S_r\longrightarrow D_r$, then $T_r\longrightarrow g_r$\\
     Four retained $x_r$ sheets in $g_r$};
  \draw[rounded corners] (4.9,1.3) rectangle (6.7,3.5);
  \draw[rounded corners] (7.3,1.3) rectangle (9.1,3.5);
  \node at (5.8,3.15) {$D_r^{+}$};
  \node at (8.2,3.15) {$D_r^{-}$};
  \draw (5.8,2.5) ellipse (.55 and .22);
  \draw (5.8,1.8) ellipse (.55 and .22);
  \draw (8.2,2.5) ellipse (.55 and .22);
  \draw (8.2,1.8) ellipse (.55 and .22);
  \node at (5.8,2.5) {$+$};
  \node at (5.8,1.8) {$-$};
  \node at (8.2,2.5) {$-$};
  \node at (8.2,1.8) {$+$};
  \node[align=center] at (7,.6)
    {Pair sheets within each $D_r$ copy.\\
     Opposite signs have equal Wall labels.};
  \node[align=center] at (7,-.65)
    {The unused $y_r$ cap remains in $M$\\
     and kills the label-comparison loop.};
\end{tikzpicture}

\caption{Attachment tree for one relation and the two ordered
compressions.  The left panel records attachments; the right panel
separates the two outer contracting-disc copies and their inner cap pairs.
Signs indicate the opposite induced orientations in each pair.
Opposite end cap types may cross in the balls $P_e$.  The unused $y_r$ cap
remains in the manifold and will kill the comparison loop in
Lemma~\ref{geo:labels}.}\label{geo:fig-stages}
\end{figure}

\begin{lemma}[Embedded framed dual candidates]\label{geo:dual-candidates}
Each $g_r$ is individually embedded and has trivial normal bundle.
It follows that
\[
 \lambda(g_r,g_r)=0,\qquad \widetilde\mu(g_r)=0.
\]
The compression leaves the single product-torus intersection
with $f_r$ unchanged.  Every other intersection involving an
$f_r$ and a $g_s$, or two distinct spheres, occurs on the cap
sheets in \eqref{geo:product-sheets}.
\end{lemma}

\begin{proof}
The chosen angular normal field extends over the embedded disc
$D_r$, so a narrow interval thickening of that disc replaces the
annulus on $T_r$ by its two disjoint parallel boundary discs.
Away from the attaching curve the thickening misses $T_r$ if
it is sufficiently small.  All cap sheets in this one sphere
come from its own $x_r$ cap and hence have the same end type.
The disjoint-body construction and the local surgery model give
embeddedness.  The phase containing the shallow-cap intersection
is disjoint from the annulus being replaced, so that point remains.
The separation choices in Lemma~\ref{geo:stages} account for all
other possible intersections.

The interval thickening is an oriented three-chain in $M$ whose
boundary is the difference between the original annulus and the
two replacement discs.  Consequently $[g_r]=[T_r]$ in
$H_2(M;\mathbb Z)$, up to a harmless orientation sign.  The torus
has an explicit trivial normal bundle, with one vector lateral
to $\beta_r$ in the base and one vector radial in the normal
disc to the base.  Thus its ordinary self-intersection is zero.
Homology invariance of that pairing gives $g_r\cdot g_r=0$.
For an embedded oriented sphere this is the Euler number of its
oriented normal plane bundle.  Such a bundle on $S^2$ is classified
by its Euler number, so $g_r$ has a disjoint normal pushoff.
The pushoff computes zero in the full group-ring diagonal pairing.
Embeddedness gives no double points at all in the Wall invariant,
and in particular gives the stated reduced invariant.
\end{proof}

Opposite intersection signs do not suffice for the off-diagonal
calculations: their labels must agree in $\pi_1(M)$, not merely
in $\pi_1(G)$.  The unused cap in each upper stage supplies this
extra assertion.

\begin{lemma}[Cancellation in the exterior]\label{geo:labels}
Within one contracting-disc copy on a sphere, pair the two
oppositely oriented $x$-cap sheets at their crossings with any
fixed opposite-type sheet.  The paired intersections have equal
Wall labels in $\pi_1(M)$ and opposite signs.  All cap-sheet
contributions to $\lambda(f_r,g_s)$ and to
$\lambda(g_r,g_s)$ for $r\ne s$ therefore cancel.
\end{lemma}

\begin{proof}
Use the product attachment collars to make the carrier and its
comparison homotopy explicit in the exterior.  Choose a common small
radius $\varepsilon$ for the deleted base tubes, smaller than all the
finitely many collar lengths.  Take $T_s$ at radius $2\varepsilon$.
Shorten the initial collar of $S_s$ at radius $3\varepsilon$ and join
it back to $T_s$ by the radial annulus between these two radii; call
this retained upper stage $S_s$ as well.  This collar shortening
changes neither its two basis curves nor its fundamental group.
All caps are retained.  Choose subsequent parallel displacements
smaller than $\varepsilon/10$ and than the positive clearances from
the other deleted tubes.  The cap-parallel collapse and the band
comparison below can be fixed on this radial collar, so their images
stay outside every deleted tube.

For this retained upper stage of entry $s$ form the abstract complexes
\[
 K_x=S_s\cup_{x_s}D_x,\qquad
 K=S_s\cup_{x_s}D_x\cup_{y_s}D_y.
\]
Attaching annuli are included in the cap notation.  The actual
configuration gives a map $j:K\longrightarrow M$ with the
collar just specified.  The caps may cross other cap images;
none of those images has been deleted from $M$.  The presentations
\begin{equation}\label{geo:carrier-groups}
 \pi_1(K_x)=\langle x_s,y_s\mid x_s\rangle
               =\langle y_s\rangle,
 \qquad
 \pi_1(K)=\langle x_s,y_s\mid x_s,y_s\rangle=1
\end{equation}
are presentations of these abstract complexes.

Collapse the small parallel displacement in the compressed disc
$D_s$ and reglue its cut $x_s$ band.  This gives a map
$c:D_s\longrightarrow K_x$ identifying corresponding points
of its two cap copies with the same point of $D_x$.  The actual
map of $D_s$ into $M$ is homotopic to $j\circ c$ by the normal
intervals of its cap parallels and the matching band collars.
The homotopy stays in a neighborhood of the upper stage, away
from every deleted base tube.  It is a homotopy of maps and is
allowed to cross the other cap images.

Fix one of the two outer copies of $D_s$ in $g_s$.  From a point
on its retained upper body choose paths $\gamma_+,\gamma_-$
to corresponding crossings on its two inner cap sheets.  In a
crossing chart these sheets and the fixed other sheet are
\[
 D^2\times\{a_+\},\quad D^2\times\{a_-\},
                 \quad \{b\}\times D^2.
\]
The short comparison path
$\tau(u)=(b,(1-u)a_++u a_-)$ is simultaneously a path on the
fixed other sheet and the normal interval between the paired
cap points.  Choose the paths on the fixed sheet using this
segment.  Up to conjugation by the common path from the sphere's
whisker, the quotient of the two Wall labels is represented by
\[
                  \gamma_+\,\tau\,\gamma_-^{-1}.
\]
Under the collapse comparison this becomes
$c(\gamma_+)c(\gamma_-)^{-1}$ in $K_x$: the segment $\tau$
collapses to the single point of $D_x$.  The loop might represent
a nonzero power of $y_s$ in $K_x$; Equation~\eqref{geo:carrier-groups}
kills it in $K$, using the actual unused $y_s$ cap.  Composing
its nullhomotopy with $j$ proves equality of the labels in $M$.
Arbitrary conjugators are already part of the attaching paths;
they do not alter the simple connectivity of $K$.

The two inner cap copies cap the two boundary components of the
annulus cut from $S_s$, so have opposite induced orientations.
Their signs against any fixed sheet are therefore opposite.
Reversing the orientation of an outer contracting-disc copy
reverses both signs and preserves this conclusion.  Apply the
pairing separately within each outer copy, and separately for
each fixed cap sheet on the other surface.  This accounts for
every cap crossing, including repeated end labels and opposite
end labels within one entry.  In particular no comparison of
paths between the two outer phases, and no contraction of a
base meridian in $M$, has been used.
\end{proof}

\begin{corollary}[The exact algebraic input]\label{geo:input-invariants}
There are orientations and whiskers for the discs and spheres such
that, for all $r,s\in\mathcal R$ and
$R=\mathbb Z[\pi_1(M)]$,
\begin{equation}\label{geo:invariants}
 \lambda(f_r,g_s)=\delta_{rs}\,1\in R,
 \qquad \lambda(g_r,g_s)=0\in R,
 \qquad \widetilde\mu(g_r)=0.
\end{equation}
In particular the normal second Stiefel--Whitney number of every
$g_r$ is zero.  These data satisfy all algebraic and boundary
hypotheses of the unrestricted disc-embedding assertion.
\end{corollary}

\begin{proof}
Lemmas~\ref{geo:dual-candidates} and~\ref{geo:labels} leave exactly
one signed group element in $\lambda(f_r,g_r)$ and no other
contributions.  Orient $g_r$ and change its whisker so that this
coefficient is $+1$.  Whisker changes multiply or conjugate zero
entries by units and hence preserve the other vanishings.
Triviality of the normal bundles proves the Stiefel--Whitney
assertion.  The manifold is oriented, so its orientation character
is trivial and the involution on $R$ is $h\mapsto h^{-1}$.
The reduced Wall invariant vanishes even before imposing its
additive relations $h-h^{-1}$ and the identity coefficient.
The discs have disjoint embedded boundary circles and the proper
collar model.  There is no requirement to make their mutual
intersections or self-intersection invariants vanish.
\end{proof}

\subsection{Topological disc images give smooth cuts}

\begin{lemma}[Moving the graph map off the protected discs]
\label{geo:topological-cuts}
Suppose there are pairwise disjoint proper locally flat embeddings
$\overline f_r:D^2\longrightarrow M$ with
$\overline f_r|_{S^1}=f_r|_{S^1}$.  Then $p_0$ is homotopic,
relative to $\partial G$, to a map $p:G\longrightarrow\Gamma$
with the following properties.  The bases, the short return
annuli in their deleted tubes, and the disc images lie in
$p^{-1}(\Gamma_h)$.  For one interior point of each $q_e$ edge,
the inverse image $J_e$ is a compact smooth oriented
three-manifold with exactly one boundary torus.  Additional
closed components of $J_e$ are allowed.
\end{lemma}

\begin{proof}
Join the boundary of each actual embedded disc back to
$\alpha_r$ by the radial annulus in the deleted base tube.
The bases and these return annuli lie in $V$ and hence already
map into $\Gamma_h$.  The boundary map of each disc lies in
$\Gamma_h$ and is nullhomotopic in $\Gamma$ because the actual
disc gives a filling.  The inclusion
$\pi_1(\Gamma_h)\longrightarrow\pi_1(\Gamma)$ is injective,
so it has a filling in $\Gamma_h$.  The original filling and
this new filling are homotopic relative boundary since
$\pi_2(\Gamma)=0$.  Transfer this homotopy to the disc image
using the embedding parametrization.  This step uses no
homotopy between an output disc and its input cap.

Let $B$ be the union of the bases, return annuli, and disc
images, and put $A=B\cup\partial G$.  Prescribe the just
constructed homotopies on the discs and keep the other parts
fixed.  These prescriptions agree on intersections: disc
boundaries are fixed, disc interiors are mutually disjoint and
proper in $M$, and the only intersections with the returned
pieces are the stated attaching circles.  Pasting finitely
many compact closed sets gives a continuous homotopy on $A$.
Only $B$ is to map into $\Gamma_h$; the $q$-collar portions
of $\partial G$ keep their prescribed product values.

Here the extension to $G$ need not assume that the topological
disc images are smooth or that their inclusion is a cofibration.
A finite graph is an absolute neighborhood retract for metric
spaces.  Apply its neighborhood-extension property to the
prescribed map on the closed subset
\[
       (G\times\{0\})\cup(A\times[0,1])\ \subset G\times[0,1].
\]
It extends to an open neighborhood $O$ of that subset.  Compactness
of $A\times[0,1]$ gives a neighborhood $U$ of $A$ with
$U\times[0,1]\subset O$.  Choose a continuous function
$\chi:G\longrightarrow[0,1]$ equal to one on $A$ and supported
in $U$.  If $\widetilde H$ is the neighborhood extension, then
\[
             H(x,t)=\widetilde H(x,\chi(x)t)
\]
is defined on all of $G\times[0,1]$: when $\chi(x)>0$ it uses
$U\times[0,1]$, and otherwise it uses the time-zero slice.
It is the required extension, fixed on $\partial G$.  Its final
map sends $B$ into $\Gamma_h$.

Choose mutually disjoint small open intervals about the midpoints
of the $q_e$ edges.  Their closed smaller intervals miss
$\Gamma_h$, so their inverse images miss the protected compact
set $B$.  On $\partial G$ the map is already the fixed coordinate
on the single collar $C_e=T^2\times[0,1]$.  In a small boundary
collar straighten its real interval coordinate to be constant
in the inward normal direction, using an interpolation fixed
on the boundary.  Use a slightly larger target interval to
localize this operation away from $B$.  Relative smooth
approximation of real functions, followed by relative
transversality in the remaining interior, makes the map smooth
and transverse near a smaller middle level, preserving the
boundary product and staying off $B$.  The intervals for
different edges have disjoint supports.  These changes are
homotopies within intervals, so the resulting $p$ is still
homotopic to $p_0$ relative boundary. The regular-value and boundary
regular-fiber statements are those of \cite[Section~2, pp.~10--13]{Milnor1965};
the fixed product collar supplies regularity on the boundary.

Each level $J_e$ is therefore a compact smooth hypersurface,
cooriented by the interval and oriented by $G$.  Its boundary
is exactly the boundary level $T^2\times\{\mathrm{point}\}$
in $C_e$; $Q$ contributes no boundary point at that level.
There is one component with this boundary torus, and every
other component is closed.  All the topological disc images
and returned pieces remain disjoint from the cuts.
\end{proof}

For the deformation argument, fix the following reference representation
of the $h$-subgroup, with $\mathfrak l=\mathfrak{sl}_2(\mathbb C)$:
\begin{equation}\label{geo:matrices}
 \rho(h_0)=\begin{pmatrix}2&0\\0&1/2\end{pmatrix},\qquad
 \rho(h_+)=\begin{pmatrix}1&1\\0&1\end{pmatrix},\qquad
 \rho(h_-)=\begin{pmatrix}1&0\\1&1\end{pmatrix}.
\end{equation}
The length-five word set $\mathcal H$ was chosen so that the operators
$\operatorname{Ad}\rho(c)$ for $c\in\mathcal H$ span
$\operatorname{End}(\mathfrak l)$; the calculation is given
at its point of use in Lemma~\ref{geo:spanning}.

\begin{proposition}[The geometric input and the forced cut data]
\label{geo:construction}
For every $d\geq5$, with the finite word set $\mathcal H$ defined above,
the preceding construction produces
a compact connected oriented smooth four-manifold $M$ with
nonempty boundary, a finite family $F=(f_r)_{r\in\mathcal R}$
of smooth proper disc maps with disjoint embedded boundary,
and individually embedded framed sphere maps
$(g_r)_{r\in\mathcal R}$ satisfying \eqref{geo:invariants}.

If the circles $\partial F$ bound pairwise disjoint proper
locally flat embedded discs in $M$, the following data exist:
\begin{enumerate}
\item The compact oriented smooth four-manifold $G$ has a
homotopy equivalence $p:G\longrightarrow\Gamma$, where
$\Gamma$ is the rose on $h_0,h_+,h_-,q_1,\ldots,q_d$.
The boundary is $Q\cup\bigcup_e C_e$, with the free group
and peripheral coordinates of Lemma~\ref{geo:graph} and the
longitude--meridian interchange \eqref{geo:hopf-collar}.
\item There are disjoint cooriented smooth cuts $J_1,\ldots,J_d$,
each a compact oriented three-manifold with one boundary torus
and possibly additional closed components.  Their boundary
tori are the fixed product cuts in the $C_e$.
\item Cutting $G$ along all $J_e$ gives a compact oriented smooth
four-manifold $X$, possibly disconnected.  With half collars
absorbed into $Q$, its boundary is obtained by gluing $Q$ to
the two oppositely oriented copies of each $J_e$ along their
boundary tori.  Every torus-bearing wall component has just
that one torus; closed wall components remain separate.
Regluing the paired cuts recovers $G$.
\item The map on $X$ lands in the graph cut open at the chosen
$q_e$ points, namely $\Gamma_h$ with the resulting half-edges;
it is constant on every copy of a cut at its endpoint value.
All words \eqref{geo:relations} are trivial in the fundamental
group of the component $X_Q$ containing $Q$, with the original
basings in $Q$.
\item Extend \eqref{geo:matrices} to $\pi_1(\Gamma)$ by
$\rho(q_e)=1$.  The associated adjoint local system on $G$
restricts to $\rho$ on the $h$ generators, to the identity on
all $z_a$, and to a trivial system on each $J_e$, with matching
reference frames on its two copies induced by the graph point.
\end{enumerate}
The implication uses only the asserted embedded discs.  It imposes
neither output-dual conditions nor an additional framing or relative
homotopy condition on those discs.
\end{proposition}

\begin{proof}
The input assertions are Corollary~\ref{geo:input-invariants}.
Use Lemma~\ref{geo:topological-cuts} to obtain $p$ and $J_e$.
Since $p$ is homotopic to $p_0$, Lemma~\ref{geo:graph} makes
it a homotopy equivalence.  Cutting along the disjoint level
hypersurfaces gives the stated collared manifold $X$, whose
corners can be rounded.  Connectedness of $X$ was not used
and is not asserted.  The connected region $Q$ singles out
one component $X_Q$.

For each $r$, the union of $B_r$, its return annulus and its
new disc misses every cut and meets $Q$.  It therefore lies
in $X_Q$.  In that union the curve $\alpha_r$ bounds the
returned disc, so the base boundary, which represents
$[\alpha_r,\beta_r]$, is nullhomotopic.  Its word in $Q$ is
exactly $R_r$, up to the already fixed orientation and basing
conventions.  Any overall conjugate or inverse has the same
nullity, so all based relations in \eqref{geo:relations}
hold in $\pi_1(X_Q)$.  This group argument does not assert
that an entire base has been compressed to an embedded disc.
The map to the cut graph and all the peripheral assertions
follow directly from the fixed product boundary values.
Pullback from the graph gives the last local-system assertion.
\end{proof}

The remaining sections rule out the cut data in
Proposition~\ref{geo:construction}.  Their smoothness comes from
regular levels of the ambient graph map away from the protected disc
images; the hypothetical discs themselves have not been smoothed.

\newpage
